\documentclass[11pt]{article}
\usepackage{amsmath,amssymb,amsthm,geometry}\geometry{margin=1in}
\newtheorem{theorem}{Theorem}\newtheorem{proposition}[theorem]{Proposition}
\title{Whole-scope block decomposition with passive world responses}
\date{16 November 2026}
\begin{document}\maketitle
\paragraph{Audit and target.} Independently reviewed theorem at the stated finite whole-scope interface. The result constructs an exact native block decomposition, including empty passive blocks, rather than factoring a supplied joint into independent child laws. It supplies a specific ingredient for root/child feedback assembly. It does not assert an initial algebra of arbitrary shared-world systems. Scope-history writer grades and terminal invocation grades remain distinct.

\section{Independent finite boundary interface}
Fix one prescope configuration at common world $w$, a finite participant set $I$, and a finite declared boundary carrier $U$ of actual world arrows out of $w$. Each participant has a finite raw row $(Z_i,\mu_i)$ and a boundary function $d_i:Z_i\to U$. All zero atoms are retained. Let
\[
 Z_I=\{(u,(z_i)_{i\in I}):d_i(z_i)=u\text{ for all }i\}.
\]
This notation keeps the boundary coordinate explicitly. A supplied normalized law $q$ on $Z_I$ has participant marginals $\mu_i$ and boundary marginal $\nu$. For $I\ne\varnothing$, the boundary coordinate is determined by any participant. No law on $U$ is selected independently of $q$.

At boundary $u$, supply an augmented coefficient $L_u$ and coefficient legs $K_{i,z_i}\to L_u$. Work in an admitted category where the required finite limits in the slice over $L_u$ exist. Define the actual joint coefficient
\[
 K_{I,z}=\prod_{i\in I}^{/L_u}K_{i,z_i}.
\]
The empty product in this slice is $L_u\xrightarrow{1}L_u$, not generally the ambient monoidal unit. This paragraph is an actual limit doctrine, not an arbitrary supplied coefficient cone.

For full scope content independently supply a boundary family $B$, families $Q_i\to B$, and their actual full-family limits $Q_J=\prod_{i\in J}^{/B}Q_i$. Such limits include shapes never realized and actual coefficient objects/maps. Explicitly supply a boundary body $\beta:U\to|B|$ and participant bodies $b_i:Z_i\to|Q_i|$ such that $|\pi_i|b_i=\beta d_i$ at every raw atom, including zeros. These induce a block body on every matching block tuple, whether or not it is attained by a global tuple, and a joint body into $Q_I$. The empty block body is precisely $\beta$. The empty full family is $B\to B$. This scope content is not identified with the tensor of the actors' completed terminal grades.

\paragraph{Boundary factorization is not automatic.}
Take one participant, $U=\{*\}$, $Z_1=\{0,1\}$ with fair law, $Q_1=B=\{0,1\}$ and grade projection identity, with raw body $b_1(z)=z$. A flat one-participant body exists, but no $\beta(*)$ equals both bits. An empty block with raw carrier $U$ and full family $B$ therefore has no compatible body. Raw and coefficient regrouping alone would not establish a native full-body decomposition.

An alternative declared interface refines the raw boundary to a carrier $\widetilde U$ with projection $p:\widetilde U\to U$ and body $\widetilde\beta:\widetilde U\to|B|$. Supply lifts $\widetilde d_i:Z_i\to\widetilde U$ satisfying $p\widetilde d_i=d_i$ and $|\pi_i|b_i=\widetilde\beta\widetilde d_i$ at all atoms. Matching is then on the whole refined boundary, not just its world-arrow projection. The rebase uses $p(\widetilde u)$, and the scope law must be supplied on this refined full matching carrier with its exact marginals. Scope coefficients are actual limits over declared $L_{\widetilde u}$ with supplied participant coefficient legs; the full writer boundary remains the entire $B$, including unused shapes. Empty responses retain $\widetilde u$, its marginal law, coefficient $L_{\widetilde u}$ and body $\widetilde\beta(\widetilde u)$. All boundary maps in comparisons must preserve these world, content and coefficient fields.

One possible refinement records the actual common writer-boundary witness together with $u$, but its availability and all-raw compatibility are extra data. This may change the matching carrier: tuples whose world arrows agree but writer-boundary witnesses differ are excluded. It is not a conservative replacement of a previously supplied world-only joint unless that law and every retained raw tuple already admit the required refinement. No support deletion is used to repair this defect. The following theorem uses the first, explicitly factorized branch; it also applies verbatim to a fully supplied refined branch after replacing its boundary data.

\section{Restriction and exact reassembly}
Partition $I$ into disjoint blocks $I=\coprod_{j\in J}I_j$ indexed by a finite nonempty set $J$; allow empty blocks. Define $Z_j$ by the same matching formula on $I_j$. Thus $Z_j=U$ for an empty block. Restriction $p_j:Z_I\to Z_j$ preserves the boundary coordinate; let $q_j=(p_j)_*q$. Each block has boundary marginal $\nu$. Write
\[
 Z_{\rm blk}=\prod_{j\in J}^{/U}Z_j,
 \qquad \Phi:Z_I\longrightarrow Z_{\rm blk}
\]
for the simultaneous restriction map, and retain the entire law $\lambda=\Phi_*q$ on this matching block carrier. Its marginals are $q_j$. The data of the $q_j$ alone do not determine $\lambda$.

\begin{theorem}[Full native block decomposition]
The map $\Phi$ is a bijection on all raw atoms. It identifies $q$ with the retained $\lambda$, joint context coefficients with the iterated slice limits of block coefficients, and the complete scope content family with the corresponding iterated full-family limit. These are actual coefficient and full-grade isomorphisms, natural for compatible primitive diagram maps and exact pushed joint laws. All iterated regroupings preserve the same flat participant coordinates, full supplied law and limit projections, so their coherence diagrams commute.
\end{theorem}
\begin{proof}
A matching family of block tuples has one common $u$ and one $z_i$ for every participant, since the blocks partition $I$. Concatenating their coordinates is inverse to $\Phi$. Empty blocks contain only that same forced $u$, so introduce no independent coordinate. This holds for zero atoms and for boundary values of zero marginal mass. Law preservation is the definition of $\lambda$, not an independence calculation.

For coefficients, a cone into all block products over $L_u$ is exactly a cone into all primitive $K_{i,z_i}$ over $L_u$. Empty block products impose only the already required map to $L_u$. The two universal properties give inverse actual augmented arrows. Their composites equal identity because all primitive projections and the retained boundary projection agree; when the participant set is empty, this is the terminal-object uniqueness in the slice. The same argument in the full-family slice over $B$ yields the full-grade isomorphism, including every unused shape and coefficient arrow. Induced raw bodies commute with it because their primitive projections agree. Thus this is an actual native row/grade isomorphism.

A map of primitive diagrams and boundaries induces the raw tuple map and actual limit arrows. Both routes through $\Phi$ have identical primitive coordinates and coefficient/full-family projections, proving naturality. Exact pushforward of the global law gives that of every block marginal by composition of functions. For three or more regroupings, all routes flatten to the same primitive tuple and the same limit cone, proving coherence without asserting independence or an inverse tensor--limit comparison.
\end{proof}

\paragraph{Admission is an additional law.}
If the interface allows every normalized law on the full matching carrier with prescribed component marginals, all block marginals and the retained $\lambda$ are admitted by the displayed equations. Conversely any admitted block rows with common boundary marginal and any admitted coupling $\lambda$ on their matching product reassemble by $\Phi^{-1}$ to a flat law with the required participant marginals. For a restricted admission policy this must instead be proved in both directions. Selected conditional-independent gluing, a fixed coupling chooser, causal restrictions or a prescribed set of legal child scopes can exclude a block marginal or reassembly. The theorem is a native representation result on admitted data, not an assertion that all physically desired scope menus satisfy the closure law.

\section{Passive blocks are responses, not extra sampled events}
For an empty block, the raw response is $u\in U$ with law $\nu$, coefficient $L_u$ as the neutral slice object, and full boundary content $B$. Its successor is the passive rebase $R_u$ of every old vertex in that block. The raw $u$ is the same global world arrow used by active blocks. It must not be freshly and independently sampled. When working fibrewise at a specified $u$, the passive response is deterministic; writing its marginal row above merely packages all these fibres.

An empty block need not be a legal intrinsic command of the child system, and no extra paid event is introduced. It is an external response derived from the child's supplied rebase action. A construction demanding that every block be represented by a nonempty child scope therefore fails on a root-only active scope. The repair is the explicit passive-response branch, not a declaration that the child suddenly has an enabled event. If the boundary coefficient is the augmentation unit, its neutral slice object is that unit; for general $L_u$ it is not. Inserting an ambient unit would require an additional arrow to $L_u$ and need not be neutral or even available.

For a concrete boundary test, take two possible world arrows with positive global probabilities, a root choosing between them and a child with no legal active command. Passive rebase changes the child's world-bearing local state. The global root-only scope is well defined. Its child restriction cannot be a selected intrinsic child command, but is exactly the boundary-driven passive response. A singleton raw block forgetting $u$ cannot distinguish the two passive successors.

\section{What this proves for a root constructor}
Partition a finite invocation configuration's primitive vertices into the root block and immediate-child subtrees. Assume: old addresses/ports survive active steps and passive rebase; completed marker and completed local terminal body are rebase-stable; active successors and passive responses use the same actual world arrow; births initialize at the actual post-world with finite all-raw initialization. If initialization is random, every raw initializer witness, its actual coefficient factor and its joint law with the paid outcome must already be retained in the declared scope data on both sides. Regrouping performs no new draw and uses the supplied initialization comparison at the matched post-world, not an equation equating rebased seeds with fresh seeds. Under block admission closure and the limit doctrine above, root splitting and grafting preserve every whole-scope row, law, context coefficient and scope writer grade by the theorem. They preserve post-states by checking active vertices, passive vertices and births individually. The same flat occurrence set and global scope command remain present.

This gives a native constructor isomorphism \emph{if} both sides have independently been defined with precisely these shared-world, whole-scope operations. It does not infer those operations from ordinary independent open composition. Nor does it supply a well-defined endofunctor on arbitrary opaque systems merely by displaying a configuration bijection.

\paragraph{Candidate functor on rigid flat-diagram maps.}
At the retained operational level, a map that supplies compatible primitive occurrence maps, world/rebase maps, command matching, complete raw maps and actual coefficient/full-grade maps acts on each global scope by the induced flat diagram map, transports $q$ to its pushed target command, and uses the unique limit arrows. Rebracketing by root blocks commutes with this action. Identities and composition follow from tuple functions, exact pushforward and limit uniqueness. Thus block decomposition is functorial on this explicitly rigid diagram-map category. Maps that merge participant domains, replace an opaque node with an expanded network, or provide only marginal effects are not included without additional structure.

\section{Two residuals before native initiality}
\paragraph{Scope writer versus terminal invocation grade.}
The theorem concerns the limit-built scope content family. A terminal invocation grade may instead be the entire formal polynomial $R_UC$, with content tensor of all completed actor factors. These are different constructions. Even at the shape level, taking $Q_1=Q_2=B$ with both maps identity gives a joint family with $|B|$ shapes, while the unconstrained product has $|B|^2$. For a two-element $B$ these cardinalities differ. No full-grade isomorphism is supplied by equality of realized terminal observations. An initiality theorem retaining both scope history and terminal content must specify their combined target grade and all full-body maps, or keep them as separate explicitly typed fields with their compatibility theorem. Passive rebase stability of completed terminal bodies addresses only one of these obligations.

\paragraph{Functor domain and fold uniqueness.}
The rigid flat-diagram functoriality does not yet define an endofunctor $T_P$ on a category broad enough to contain arbitrary target algebras. For such an algebra $\alpha:T_P(X)\to X$, its map must preserve or explicitly transport the whole participant diagram and global coupling, not merely primitive marginals. A recursive fold must land in that same category and have smaller finite invocation height at each recursive child call. It must be determined on menus, whole-scope raw and coefficient maps, scope writer families and terminal formal grades. Actual limits remove hidden coefficient-map freedom only over specified primitive diagram maps; they do not determine an unspecified occurrence decomposition or an opaque target's control interface.

The finite-tree constructor can therefore be natively coherent on the rigid presented subcategory while a general algebra initiality claim remains unproved. A completed theorem needs either a precise algebra category in which all target arrows retain the required decomposition, or an additional refinement/assembly doctrine proving that opaque folds canonically expose it. Neither a scalar trace law nor the ordinary set-theoretic fold provides that missing field.

\paragraph{Current advance.} Entire supplied joint laws, actual slice-limit coefficients and full scope families do admit coherent root/child regrouping, with canonical passive empty blocks. The decisive correction to naive decomposition is to retain global coupling and boundary-driven passive response, not to independently execute children. This resolves the flat block representation issue and isolates the remaining algebra-domain and terminal-grade questions. It is not a claim that all feedback obligations are complete.
\end{document}
